\section{计算机网络体系结构}

\subsection{计算机网络的概念与分类}

\begin{mydef}{计算机网络}{cn-def}
计算机网络是将分布在不同地理位置、具有独立功能的计算机系统，通过通信设备和线路连接起来，以网络协议实现资源共享和数据通信的系统。
\end{mydef}

\begin{conclusion}{计算机网络的分类}{cn-class}
\begin{itemize}
    \item \textbf{按覆盖范围}：广域网 WAN、城域网 MAN、局域网 LAN、个域网 PAN。
    \item \textbf{按传输技术}：广播式网络 vs 点对点网络。
    \item \textbf{按拓扑结构}：星形、总线形、环形、树形、网状。
    \item \textbf{按交换技术}：电路交换、报文交换、分组交换。
\end{itemize}
\textbf{408 中主要使用按覆盖范围的分类。}
\end{conclusion}

\subsection{电路交换、报文交换与分组交换}

\begin{conclusion}{三种交换技术的对比}{switching}
\begin{tablebox}{交换技术对比}
\centering
\begin{tabular}{|l|c|c|c|}
\hline
 & 电路交换 & 报文交换 & 分组交换 \\
\hline
建立连接 & 需要 & 不需要 & 不需要（虚电路需要）\\
传输单位 & 比特流 & 整个报文 & 分组（包）\\
存储转发 & 否 & 是 & 是 \\
时延 & 最小（仅传播时延）& 最大 & 中等 \\
灵活性 & 差（专用线路）& 好 & 好 \\
典型应用 & 传统电话 & 电报 & 互联网（IP）\\
\hline
\end{tabular}
\end{tablebox}

\textbf{分组交换的优点：} (1) 线路利用率高（统计复用）；(2) 不同速率/格式的终端可互通；(3) 优先级。

\textbf{分组交换的缺点：} (1) 额外开销（首部）；(2) 存储转发引入排队时延，不适合实时通信。

\textbf{408 常考：}给定场景，判断应使用哪种交换方式。
\end{conclusion}

\subsection{计算机网络的性能指标}

\begin{formula}{网络性能指标}{cn-metrics}
\begin{itemize}
    \item \textbf{速率（Data Rate）}：每秒传输的比特数（bps）。$1\text{K}=10^3$，$1\text{M}=10^6$（通信领域，非 $2^{10}$）。
    \item \textbf{带宽（Bandwidth）}：信道允许通过的频率范围（Hz），或网络通信线路所能传送数据的能力（bps）。
    \item \textbf{吞吐量（Throughput）}：单位时间内实际通过网络的数据量。
    \item \textbf{时延（Delay/Latency）}：$d_{\text{总}} = d_{\text{发送}} + d_{\text{传播}} + d_{\text{处理}} + d_{\text{排队}}$
    \begin{itemize}
        \item \textbf{发送时延}：$\dfrac{\text{数据帧长度}}{\text{发送速率}}$（与传输距离无关）
        \item \textbf{传播时延}：$\dfrac{\text{信道长度}}{\text{电磁波传播速率}}$（与发送速率无关）
    \end{itemize}
    \item \textbf{时延带宽积}：$d_{\text{传播}} \times \text{带宽}$ = 链路中最多容纳的比特数。
    \item \textbf{往返时间 RTT}：数据从发送端发出到收到确认的时间。
    \item \textbf{利用率}：信道利用率 = 有数据通过时间 / 总时间。利用率趋近 1 时时延急剧增大（排队论）。
\end{itemize}
\end{formula}

\begin{attention}
\textbf{408 高频计算：发送时延 vs 传播时延。} 发送时延 = 数据量/带宽（瓶颈在发送速率）；传播时延 = 距离/光速（瓶颈在物理距离）。两者影响的变因完全不同，选择题经常混合考察。
\end{attention}

\subsection{OSI 参考模型（7层）}

\begin{mydef}{OSI 七层模型}{osi-model}
国际标准化组织（ISO）提出的开放系统互连参考模型：
\[
\text{应用层} \to \text{表示层} \to \text{会话层} \to \text{传输层} \to \text{网络层} \to \text{数据链路层} \to \text{物理层}
\]
记忆口诀：\textbf{All People Seem To Need Data Processing}（从高到低）。
\end{mydef}

\begin{tablebox}{OSI 各层功能}
\centering
\begin{tabular}{|c|l|l|}
\hline
层次 & 功能 & 数据单元 \\
\hline
物理层 & 比特传输，定义机械/电气/功能/规程特性 & 比特（bit）\\
数据链路层 & 相邻结点间可靠传输，成帧、差错控制、流量控制 & 帧（Frame）\\
网络层 & 路由选择、拥塞控制、网际互连 & 分组/数据报（Packet）\\
传输层 & 端到端可靠/不可靠传输，复用与分用 & 报文段/用户数据报 \\
会话层 & 会话管理、同步 & — \\
表示层 & 数据格式转换、加密解密、压缩 & — \\
应用层 & 为应用程序提供网络服务 & 报文（Message）\\
\hline
\end{tabular}
\end{tablebox}

\subsection{TCP/IP 模型（4层）}

\begin{conclusion}{TCP/IP 四层模型}{tcpip-model}
TCP/IP 是事实上的工业标准，而非法律上的国际标准（OSI 是法律标准但未流行）。

\begin{tablebox}{OSI 与 TCP/IP 对应关系}
\centering
\begin{tabular}{|c|c|l|}
\hline
OSI & TCP/IP & 典型协议 \\
\hline
应用层、表示层、会话层 & 应用层 & HTTP, DNS, SMTP, FTP \\
传输层 & 传输层 & TCP, UDP \\
网络层 & 网际层 & IP, ICMP, ARP, OSPF, BGP \\
数据链路层、物理层 & 网络接口层 & Ethernet, PPP \\
\hline
\end{tabular}
\end{tablebox}

\textbf{五层参考模型（教学用，408采用）：} 物理层 $\to$ 数据链路层 $\to$ 网络层 $\to$ 传输层 $\to$ 应用层。
\end{conclusion}

\subsection{协议与服务}

\begin{mydef}{协议三要素}{protocol-triple}
网络协议（Protocol）是为进行数据交换而建立的规则，由三个要素组成：
\begin{enumerate}
    \item \textbf{语法（Syntax）}：数据与控制信息的结构或格式；
    \item \textbf{语义（Semantics）}：每种控制信息的含义以及应完成的动作；
    \item \textbf{同步/时序（Timing）}：事件实现顺序的详细说明。
\end{enumerate}
\end{mydef}

\begin{conclusion}{协议与服务的关系}{protocol-service}
\begin{itemize}
    \item \textbf{协议是「水平的」}：控制两个对等实体之间的通信规则。
    \item \textbf{服务是「垂直的」}：下层通过层间接口向上层提供功能。
    \item 第 $n$ 层实体向第 $n+1$ 层提供服务（Service），使用的对等协议在第 $n$ 层。
    \item \textbf{SAP（服务访问点）}：相邻层间交换信息的逻辑接口。数据链路层的 SAP 是 MAC 地址，网络层是 IP 地址，传输层是端口号。
\end{itemize}
\end{conclusion}

\subsection{408 高频考点速查}

\begin{conclusion}{体系结构 — 408常考点}{cn1-keypoints}
\begin{enumerate}
    \item \textbf{发送时延 vs 传播时延}：计算公式、影响因素。
    \item \textbf{OSI 7层 vs TCP/IP 4层 vs 五层模型}：各层名称、功能、典型协议、数据单元名称。
    \item \textbf{协议三要素}：语法、语义、时序。
    \item \textbf{分组交换 vs 电路交换}：各自时延构成、适用场景。
    \item \textbf{各层中间设备}：物理层（中继器/集线器）、数据链路层（网桥/交换机）、网络层（路由器）、传输层及以上（网关）。
\end{enumerate}
\end{conclusion}

\newpage
\section{习题}

\subsection{基础过关题}

\begin{enumerate}

\item \textbf{（2009年·选择题·第33题）} 在OSI参考模型中，自下而上第一个提供端到端服务的层次是
\begin{center}
\begin{tabular}{ll}
(A) 数据链路层 & (B) 传输层 \\[6pt]
(C) 会话层 & (D) 应用层
\end{tabular}
\end{center}

\ifshowsolutions\par\noindent\textbf{解析：} 传输层是第一个端到端的层次——网络层及以下都是逐跳（hop-by-hop）。\boxed{\textbf{答案：}(B)}\fi

\vspace{8pt}

\item \textbf{（2010年·选择题·第33题）} 下列选项中，不属于网络体系结构所描述的内容的是
\begin{center}
\begin{tabular}{ll}
(A) 网络的层次 & (B) 每一层使用的协议 \\[6pt]
(C) 协议的内部实现细节 & (D) 每一层必须完成的功能
\end{tabular}
\end{center}

\ifshowsolutions\par\noindent\textbf{解析：} 网络体系结构只定义层次划分和各层功能/协议接口，不规定协议的内部实现细节。\boxed{\textbf{答案：}(C)}\fi

\vspace{8pt}

\item \textbf{（2012年·选择题·第33题）} TCP/IP参考模型的网络层（网际层）提供的是
\begin{center}
\begin{tabular}{ll}
(A) 无连接不可靠的数据报服务 & (B) 无连接可靠的数据报服务 \\[6pt]
(C) 面向连接不可靠的虚电路服务 & (D) 面向连接可靠的虚电路服务
\end{tabular}
\end{center}

\ifshowsolutions\par\noindent\textbf{解析：} IP协议提供无连接、不可靠的尽力而为（best-effort）数据报服务。可靠传输由传输层TCP保证。\boxed{\textbf{答案：}(A)}\fi

\vspace{8pt}

\item \textbf{（2016年·选择题·第33题）} 假设OSI参考模型的应用层欲发送400 B的数据（无拆分）。除物理层和应用层外，其他各层在封装PDU时均引入20 B的额外开销，则应用层的数据传输效率约为
\begin{center}
\begin{tabular}{ll}
(A) 80\% & (B) 83\% \\[6pt]
(C) 87\% & (D) 91\%
\end{tabular}
\end{center}

\ifshowsolutions\par\noindent\textbf{解析：} 共5层开销（表示层、会话层、传输层、网络层、数据链路层），每层20B，共$5\times 20=100$B。总发送量$=400+100=500$B。效率$=400/500=80\%$。\boxed{\textbf{答案：}(A)}\fi

\vspace{8pt}

\item \textbf{（计算题）} 收发两端之间的传输距离为1000 km，信号在媒体上的传播速率为$2\times 10^8$ m/s。数据长度为$10^7$ bit，数据发送速率为100 kb/s。求发送时延和传播时延。

\ifshowsolutions\par\noindent\textbf{解析：}
发送时延$=10^7/10^5=100$ s。传播时延$=10^6/(2\times 10^8)=0.005$ s$=5$ ms。\fi

\end{enumerate}

\subsection{真题风格与综合训练}

\begin{enumerate}[resume]

\item \textbf{（2013年·选择题·第35题）} 主机甲通过1个路由器（存储转发方式）与主机乙互连，两段链路的数据传输速率均为10 Mbps。主机甲分别用报文交换和分组大小为10 kb的分组交换向主机乙发送1个大小为8 Mb的报文。若忽略链路传播时延、分组首部开销和拆装时间，则两种交换方式完成该报文传输所需的总时间分别为
\begin{center}
\begin{tabular}{ll}
(A) 800 ms, 1600 ms & (B) 801 ms, 1600 ms \\[6pt]
(C) 1600 ms, 800 ms & (D) 1600 ms, 801 ms
\end{tabular}
\end{center}

\ifshowsolutions\par\noindent\textbf{解析：}
报文交换：甲发完$8\text{ Mb}/10\text{ Mbps}=800$ ms；路由器再发800 ms到乙。总$=1600$ ms。

分组交换：$8\text{ Mb}/10\text{ kb}=800$个分组。甲发完所有分组$=800\times(10\text{ kb}/10\text{ Mbps})=800$ ms。第一个分组到达乙：甲发完第一个分组（1 ms）$+$路由器转发第一个分组（1 ms）$=2$ ms。最后一个分组到达乙$=800+1=801$ ms（最后一个分组甲发完时，路由器在1 ms前已完成倒数第二个分组的转发，可立即转发最后一个）。

总时间$=800+1=801$ ms。\boxed{\textbf{答案：}(D)}\fi

\vspace{8pt}

\item \textbf{（真题风格·综合题）} 画出TCP/IP四层模型与OSI七层模型的对应关系，并标注各层典型协议和数据单元名称。

\ifshowsolutions\par\noindent\textbf{答案：}
\begin{center}
\begin{tabular}{|c|c|l|l|}
\hline
OSI & TCP/IP & 典型协议 & PDU \\
\hline
应用层、表示层、会话层 & 应用层 & HTTP, DNS, SMTP & 报文 \\
传输层 & 传输层 & TCP, UDP & 报文段/数据报 \\
网络层 & 网际层 & IP, ICMP, ARP & 数据报/分组 \\
数据链路层、物理层 & 网络接口层 & Ethernet & 帧/比特 \\
\hline
\end{tabular}
\end{center}\fi

\end{enumerate}

\vspace{8pt}
\noindent\textbf{拓展阅读：}
\begin{itemize}
    \item Kurose \& Ross《计算机网络：自顶向下方法》第1章——生动直观的分组交换vs电路交换定量分析。
    \item 谢希仁《计算机网络》第1章——国内标准术语体系。
\end{itemize}
